The fastest-growing lever on a frozen LLM is no longer weights but the
\emph{harness} --- the program of prompts, tools, and control flow that
wraps the model for a task. A recent wave of systems generates this
scaffolding automatically: LLM builders propose, edit, and evolve
harness code, with reported gains over bare prompting. The implicit
premise of this ``AI4AI'' line is that when the generated code differs,
the \emph{behavior} differs --- that a population of thirty harnesses is
thirty ways of answering, not one way written thirty times.

We test that premise directly. On BIRD text-to-SQL with a frozen
mid-tier target, we take the full candidate population produced by a
standard harness-evolution loop and measure what its members actually
do, not what their source files say. Every candidate passes a
code-difference check, yet among the 12 generated candidates average
pairwise outcome disagreement is \textbf{3.8\%} and \textbf{21 of 66
candidate pairs are behaviorally identical} (28 of 78 pairs once the
shipped ReAct harness is included); the ReAct variant reproduces the
bare baseline's final SQL \textbf{byte-for-byte on every audited task}.
Execution traces resolve ``different code'' into a three-class taxonomy
--- mechanisms \emph{described} but absent from the control flow,
feedback with no causal path to the output, and genuinely multi-turn
paths that are broken. We call this \textbf{behavioral collapse}:
syntactic diversity without behavioral diversity (F1).
